\section{ACI and Blockchain Migration}

\label{sec:blockchain}
% ============================================================================

Blockchain signatures present unique challenges: transactions signed today are
verified forever. A quantum adversary can forge historical transactions if the
signing algorithm is broken.

\subsection{Lux Network Coordination}

The Lux blockchain~\cite{luxwhitepaper} is independently migrating to
post-quantum signatures. Hanzo ACI coordinates with Lux for:
\begin{itemize}[leftmargin=1.1em]
    \item \textbf{Consensus signatures}: Validator attestations migrate from
    ECDSA to ML-DSA in a coordinated network upgrade.
    \item \textbf{Transaction signatures}: New transaction types support
    ML-DSA-65 signatures alongside legacy ECDSA (secp256k1).
    \item \textbf{Address derivation}: Post-quantum addresses are derived from
    ML-DSA-65 public keys, distinct from legacy ECDSA addresses.
\end{itemize}

\subsection{Threshold Signing Migration}

Hanzo Threshold Signing~\cite{hanzothreshold} migrates from FROST (Ed25519)
and CGGMP21 (ECDSA) to post-quantum threshold schemes:
\begin{itemize}[leftmargin=1.1em]
    \item \textbf{Lattice-based threshold signatures}: Research prototypes of
    threshold ML-DSA exist but are not yet production-ready. Expected
    standardization: 2027.
    \item \textbf{Interim approach}: Hash-based threshold signatures using
    SLH-DSA with Shamir secret sharing of the hash-based signing key. This
    is less efficient but provably secure.
\end{itemize}

% ============================================================================
